\unit{Exposição 212 -- Pré-esquemas quocientes}
\sid{212}
\src{99}
\seminares{13.º ano, 1960/61, n.º 212\\Fevereiro de 1961}
\paperhead{TÉCNICAS DE CONSTRUÇÃO E TEOREMAS DE EXISTÊNCIA EM GEOMETRIA ALGÉBRICA\\
III. PRÉ-ESQUEMAS QUOCIENTES}%
  {por Alexander GROTHENDIECK}

\subsection*{Introdução}
\addcontentsline{toc}{subsection}{Introdução}

Os problemas tratados na presente exposição diferem dos considerados nas duas
anteriores, pois se procura representar certos funtores covariantes, e não mais
contravariantes, de esquemas variáveis. O procedimento de passagem ao quociente
é, contudo, essencial em muitos problemas de construção em geometria algébrica,
inclusive os das Exposições I e II ([1], [2]). Assim, a questão da
\emph{efetividade de um dado de descida} sobre um $T$-pré-esquema $X$,
relativamente a um morfismo fielmente plano e quase-compacto
$T\longrightarrow S$, equivale à questão da existência de um quociente de $X$
(que satisfaça as propriedades razoáveis examinadas adiante) pela relação de
equivalência plana em $X$ definida pelo dado de descida; as questões levantadas
em [1], A, \S~2 c serão, sem dúvida, resolvidas ao mesmo tempo que as questões
formuladas no parágrafo 2 da presente exposição. Do mesmo modo, o
\emph{esquema de Picard} de um $S$-esquema $X$ (para a definição, ver [2], C,
\S~3) pode ser definido de várias maneiras como quociente de certos outros
esquemas (de divisores efetivos, ou de imersões em um esquema projetivo) por
relações de equivalência planas, sendo tecnicamente mais simples a definição e
a construção desses esquemas auxiliares: são esquemas da forma
$\underline{\operatorname{Hom}}_S(X,Y)$ e variantes definidos em [2], C,
\S~2, cuja construção será objeto da exposição seguinte (sob hipóteses
adequadas de projetividade). Combinando, portanto, os resultados da presente
exposição e da seguinte, obtém-se a construção dos esquemas de Picard sob
hipóteses adequadas.

O problema da passagem ao quociente na categoria dos pré-esquemas ainda
apresenta várias questões não resolvidas. A mais importante é mencionada no
parágrafo 8. Atualmente, ela é o único obstáculo à construção de
\emph{esquemas de moduli sobre os inteiros para curvas de gênero arbitrário,
variedades abelianas polarizadas}, etc. Isso mostra que sua solução merece os
esforços dos especialistas em grupos algébricos.

\src{100}
\subsection*{1. Relações de equivalência e relações de equivalência efetivas}
\addcontentsline{toc}{subsection}{1. Relações de equivalência e relações de equivalência efetivas}

Sejam $\mathcal C$ uma categoria e $X$ um objeto de $\mathcal C$. Um par de
morfismos
\[
  p_1,p_2:R\rightrightarrows X
\]
é chamado de \emph{par de equivalência} em $\mathcal C$, com \emph{alvo} $X$ e
\emph{fonte} $R$, se, para todo objeto $T$ de $\mathcal C$, as duas aplicações
correspondentes
\[
  p_1(T),p_2(T):R(T)\rightrightarrows X(T)
\]
(onde, para todo objeto $Y$ de $\mathcal C$, se põe
$Y(T)=\operatorname{Hom}(T,Y)$) definem uma aplicação
\[
  R(T)\longrightarrow X(T)\times X(T)
\]
que induz uma bijeção de $R(T)$ sobre o gráfico de uma relação de equivalência
no conjunto $X(T)$. Introduz-se, entre os pares de equivalência com alvo $X$,
uma relação de equivalência evidente; uma classe de equivalência para esta
última é chamada \emph{$\mathcal C$-relação de equivalência em $X$}, ou
simplesmente uma relação de equivalência quando não houver risco de confusão.

Se $X\times X$ existe, dar uma relação de equivalência em $X$ equivale a dar um
subobjeto $R$ de $X\times X$ tal que, para todo objeto $T$ de $\mathcal C$, o
subconjunto de
\[
  (X\times X)(T)=X(T)\times X(T)
\]
correspondente a $R(T)$ seja o gráfico de uma relação de equivalência em
$X(T)$. Denotando por $p_1$ e $p_2$ os morfismos de $R$ em $X$ induzidos pelas
projeções $\operatorname{pr}_1$ e $\operatorname{pr}_2$, a condição anterior
exprime que $(p_1,p_2)$ é um par de equivalência. Também se podem exprimir
diagramaticamente em $\mathcal C$ (desde que existam $X\times X$ e o produto
fibrado $(R,p_2)\times_X(R,p_1)$) os axiomas conjuntistas de uma relação de
equivalência para $R(T)$ em $X(T)$, de acordo com o princípio geral
explicado em [2], A, \S~1. Não precisaremos disso.

Sempre que se tem um par de morfismos $(p_1,p_2)$ com a mesma fonte $R$ e o
mesmo alvo $X$, pode-se definir o \emph{conúcleo} do par como o objeto $Y$ de
$\mathcal C$ que representa o funtor covariante em $Z$
\[
  \operatorname{Hom}_{p_1,p_2}(X,Z),
\]
conjunto dos morfismos $u$ de $X$ em $Z$ tais que $up_1=up_2$. Se $Y$ existe,
fica determinado a menos de um isomorfismo único. Será denotado por
$X/(p_1,p_2)$ ou, por abuso de notação, por $X/R$, sendo esta última notação
empregada sobretudo quando

\src{101}$(p_1,p_2)$ é um par de equivalência: é então usual identificar,
na notação, a relação de equivalência definida pelo par com $R$. Observe-se
que, se $Y$ for considerado um quociente de $X$, ele depende, de fato, apenas
da relação de equivalência definida pelo par $(p_1,p_2)$.

Partamos agora de um morfismo
\[
  f:X\longrightarrow Y
\]
que permite, portanto, considerar $X$ como um «objeto sobre $Y$», e suponhamos
que exista o produto fibrado
\[
  R(f)=X\times_Y X
\]
Sejam $p_1$ e $p_2$ suas duas projeções. Então $(p_1,p_2)$ é um par de
equivalência, chamado par \emph{associado} ao morfismo $f$. Ele define, assim,
uma relação de equivalência, chamada relação associada a $f$.

Um par de morfismos $(p_1,p_2)$ com alvo $X$ e fonte comum $R$ é chamado
\emph{par de equivalência efetivo} se
\begin{enumerate}
  \item[(i)] o conúcleo $X/(p_1,p_2)=Y$ existe;
  \item[(ii)] o produto fibrado $X\times_Y X$ existe;
  \item[(iii)] o morfismo $R\longrightarrow X\times_Y X$ cujas componentes são
  $p_1$ e $p_2$ é um isomorfismo.
\end{enumerate}
Então $(p_1,p_2)$ é de fato um par de equivalência. A relação de equivalência
que ele define é também chamada \emph{relação de equivalência efetiva}.

Um morfismo $f:X\longrightarrow Y$ é chamado \emph{epimorfismo efetivo} se
\begin{enumerate}
  \item[(i)] o produto fibrado $X\times_Y X=R$ existe;
  \item[(ii)] o quociente $X/(p_1,p_2)$ existe, onde $p_1$ e $p_2$ são as duas
  projeções de $R$ em $X$;
  \item[(iii)] o morfismo $X/(p_1,p_2)\longrightarrow Y$ induzido por $f$ é um
  isomorfismo.
\end{enumerate}
Então $f$ é de fato um epimorfismo, e até um epimorfismo estrito (cf. [1], A,
\S~2.3), sendo a recíproca verdadeira se o produto fibrado $X\times_Y X$
existir. Diz-se também que o objeto quociente de $X$ definido pelo epimorfismo
$f$ é um \emph{quociente efetivo} de $X$.

As definições precedentes implicam a seguinte «correspondência de Galois»:

\src{102}
\result{PROPOSIÇÃO}{1.1} Há uma correspondência biunívoca, que preserva as
ordens naturais, entre o conjunto das relações de equivalência efetivas $R$
em $X$ e o conjunto dos quocientes efetivos $Y$ de $X$: a $R$ corresponde o
quociente efetivo $X/R$, e a $Y$ corresponde a relação de equivalência efetiva
definida pela projeção canônica $X\longrightarrow Y$ (que é definida pelo
produto fibrado $X\times_Y X$ munido de suas duas projeções).

Nas categorias muito boas (conjuntos, feixes de conjuntos etc.), todo
quociente é efetivo e toda relação de equivalência é efetiva. Isso já não
ocorre em categorias como a categoria dos pré-esquemas sobre um pré-esquema
$S$ dado, mesmo quando $S$ é o espectro de um corpo e mesmo quando nos
restringimos aos esquemas finitos sobre $S$. As questões de efetividade e até
mesmo (no caso de pré-esquemas não finitos sobre $S$) as questões relativas à
existência de quocientes revelam-se quase sempre delicadas.

\subsection*{2. Exemplo: pré-esquemas finitos sobre $S$}
\addcontentsline{toc}{subsection}{2. Exemplo: pré-esquemas finitos sobre S}

Seja $\mathcal C$ a categoria dos pré-esquemas finitos sobre $S$, supondo-se
$S$ localmente Noetheriano. Ela é equivalente à categoria oposta da categoria
dos feixes coerentes de álgebras comutativas sobre $S$, ou ainda, se $S$ é
afim com anel $A$, à categoria oposta da categoria das $A$-álgebras finitas
sobre $A$ (isto é, finitamente geradas como $A$-módulos). Conclui-se
imediatamente que existem em $\mathcal C$ os limites projetivos finitos e os
limites indutivos finitos. Isso é bem conhecido (sem qualquer hipótese de
finitude\ldots) quanto aos primeiros. Assim, o produto fibrado de pré-esquemas
$X,Y$ sobre $S$ corresponde ao produto tensorial $B\otimes_A C$ das álgebras
correspondentes; o núcleo de dois morfismos $X\rightrightarrows Y$,
correspondentes a dois homomorfismos de $A$-álgebras
$u,v:C\rightrightarrows B$, corresponde ao quociente de $B$ pelo ideal gerado
pelos elementos $u(c)-v(c)$ etc. Para os limites indutivos finitos, basta
considerar, de um lado, as somas finitas, que correspondem ao produto ordinário
de $A$-álgebras, e, de outro, os conúcleos de pares de morfismos
$X\rightrightarrows Y$, que correspondem de fato (como se verifica
imediatamente) ao subanel de $C$ formado pelos elementos nos quais os
homomorfismos $u,v:C\rightrightarrows B$ coincidem (este último é finito sobre
$A$ graças à hipótese Noetheriana). Observe-se, aliás, que se pode mostrar,
utilizando a hipótese Noetheriana, que os limites indutivos finitos e, em
particular, os quocientes assim construídos na categoria $\mathcal C$ dos
pré-esquemas finitos sobre $S$ são, na realidade, quocientes na categoria de
\emph{todos} os pré-esquemas.

\src{103}
Como dissemos em [1], há em $\mathcal C$ epimorfismos não efetivos (ou
estritos\footnote{O texto impresso omite «não» antes de «estritos». Na
categoria considerada existem produtos fibrados, e o parágrafo anterior
estabelece precisamente a equivalência entre efetividade e epimorfia estrita.};
dá no mesmo, pois existem produtos fibrados). \emph{Ignoro se as relações de
equivalência são sempre efetivas} quando não se impõe nenhuma hipótese de
planicidade. Nessa direção, obtive apenas resultados positivos muito parciais,
indispensáveis à demonstração do teorema fundamental da teoria formal de
moduli (ver [2], B, teo. 1). Assinalemos que, no problema proposto, é fácil
reduzir-se ao caso em que $S$ é o espectro de um anel artiniano local $A$, com
corpo residual algebricamente fechado. Mas, mesmo se $A$ for um corpo, a
resposta não é conhecida.

Também se pode considerar um pré-esquema $X$ sobre $S$ que já não se suponha
finito sobre $S$, mas que seja munido de uma relação de equivalência $R$ em
$X$ tal que $p_1:R\longrightarrow X$ seja um morfismo finito. Diz-se então que
$R$ é uma relação de equivalência finita. Suponhamos, para simplificar, que $S$
e $X$ sejam afins (logo, $R$ é afim, de modo que a situação se reduz a uma
questão de álgebra comutativa pura). \emph{Mesmo nesse caso, ignora-se se
existe um quociente $X/R=Y$ e se o morfismo canônico $X\longrightarrow Y$ é
finito}. (O caso mais simples é aquele em que $S$ é o espectro de um corpo $k$
e $X$ é o espectro de $k[t]$, isto é, a reta afim.) Naturalmente, se os dois
problemas precedentes tiverem soluções afirmativas, pode-se concluir, na
situação presente, que $R$ é efetiva. Observemos que o problema da
\emph{existência} de um quociente $Y$ e da finitude de
$f:X\longrightarrow Y$ se apresenta exatamente nos mesmos termos se, em vez
do grafo de uma relação de equivalência em $X$, houver apenas um pré-grafo de
equivalência em $X$, no sentido do parágrafo 4.

A questão da passagem ao quociente por uma relação de equivalência finita mais
ou menos arbitrária surge na construção de pré-esquemas por «colagem» de
pré-esquemas dados $X_i$ ao longo de certos subpré-esquemas fechados; a lei de
colagem exprime-se precisamente por uma relação de equivalência finita no
pré-esquema soma $X$ dos $X_i$. Deve-se esperar também que a solução dos
problemas aqui propostos, e de diversas variantes, seja uma condição prévia ao
desenvolvimento de uma técnica geral de construções não projetivas, na direção
iniciada em [2].

O único fato positivo geral conhecido pelo autor é o seguinte:

\result{PROPOSIÇÃO}{2.1} Sejam $S$ um pré-esquema localmente noetheriano, $s$
um ponto de $S$ e $\Omega$ uma extensão algebricamente fechada de $k(s)$.
Considere o

\src{104}«funtor fibra» correspondente $F$, que associa a todo $S$-esquema
$X$ finito sobre $S$ o conjunto dos pontos de $X_s$ com valores em $\Omega$.
Esse funtor (trivialmente exato à esquerda) é \emph{exato à direita}, isto é,
comuta com os limites indutivos finitos e, em particular, com os conúcleos de
pares de morfismos.

Usando esse resultado para todos os «pontos geométricos» de $S$, deduz-se que
a categoria «quociente» $\mathcal C'$ de $\mathcal C$, obtida raciocinando
«módulo morfismos radiciais sobrejetivos» (isto é, adjungindo formalmente
inversos aos morfismos em questão), é uma categoria «geométrica», isto é,
satisfaz as mesmas propriedades «de natureza finita» que a categoria dos
conjuntos. Em particular, toda relação de equivalência nela é efetiva. Isso
implica que, se $R$ for uma relação de equivalência em um $X$ finito sobre
$S$, então o morfismo canônico
\[
  R\longrightarrow X\times_Y X\qquad (\text{onde }Y=X/R)
\]
é \emph{radicial e sobrejetivo} (de fato, uma imersão fechada sobrejetiva,
pois é um monomorfismo).

\subsection*{3. O caso de um grupo de operadores}
\addcontentsline{toc}{subsection}{3. O caso de um grupo de operadores}

Suponhamos novamente que $\mathcal C$ seja uma categoria qualquer. Sejam
$G,X$ objetos de $\mathcal C$, e suponhamos que $G$ seja um
$\mathcal C$-grupo de operadores sobre o objeto $X$ de $\mathcal C$. Isso
significa (cf. [2], A, \S~1) que, para todo objeto $T$ de $\mathcal C$, foram
dadas uma estrutura de grupo em $G(T)$ e uma estrutura de $G(T)$-conjunto em
$X(T)$, de tal modo que, variando $T$, as estruturas em questão «variam
funtorialmente» com $T$. Se os produtos $G\times G$ e $G\times X$ existirem
em $\mathcal C$, tal estrutura também poderá ser definida por um par de
morfismos
\[
  G\times G\longrightarrow G,
  \qquad \pi:G\times X\longrightarrow X,
\]
sujeito à condição de que, para todo objeto $T$ de $\mathcal C$, as leis de
composição correspondentes para os conjuntos $G(T)$ e $X(T)$ façam de $G(T)$
um grupo que age sobre $X(T)$. A tradução desse axioma pela comutatividade de
certos diagramas em $\mathcal C$ é fácil, mas fastidiosa e, de fato,
perfeitamente inútil em todos os casos que conheço.

Suponhamos que $G\times X$ exista e consideremos os dois morfismos
\[
  p_1,p_2:G\times X\rightrightarrows X
\]
dados por

\src{105}
\[
  p_1=\operatorname{pr}_1,\qquad p_2=\pi.
\]
Verifica-se imediatamente que $(p_1,p_2)$ é um par de equivalência se, e
somente se, para todo objeto $T$ de $\mathcal C$, a aplicação
\[
  G(T)\times X(T)\simeq(G\times X)(T)
  \longrightarrow X(T)\times X(T),
\]
definida por esse par é injetiva, isto é, se o grupo $G(T)$ age
\emph{livremente} sobre o conjunto $X(T)$, ou seja, se $g\in G(T)$,
$x\in X(T)$ e $g\cdot x=x$ implicam que $g$ é o elemento identidade de
$G(T)$. Diz-se então que $G$ \emph{age livremente sobre $X$} (ou que $X$ é um
\emph{$\mathcal C$-espaço principal sob $G$}). A relação de equivalência
associada ao par $(p_1,p_2)$ chama-se então \emph{relação de equivalência
definida pelo grupo $G$ que age livremente sobre $X$}. Quando $X\times X$
também existe e se considera o morfismo
\[
  p:G\times X\longrightarrow X\times X
\]
definido pelo par $(p_1,p_2)$, a condição de que $G$ aja livremente significa
que $p$ é um \emph{monomorfismo}.

Naturalmente, mesmo quando $G$ não age livremente sobre $X$, procuram-se
critérios para a existência de um quociente de $X$ por $G$, isto é, do
conúcleo do par anterior $(p_1,p_2)$.

O conúcleo em questão será muitas vezes denotado por $X/G$, ou,
preferencialmente, por $G\backslash X$ se $G$ agir à esquerda (reservando-se a
notação anterior para o caso em que $G$ age à direita). Observe-se que, mesmo
quando existe a «imagem» de $G\times X$ por $p$ (definida, por exemplo, como o
menor subobjeto de $X\times X$ pelo qual $p$ se fatora), digamos $R$, ela quase
nunca é uma relação de equivalência em $X$. Se então se tenta passar
diretamente ao quociente por $R$ (ou, mais precisamente, pelo par de morfismos
de $R$ em $X$ induzidos pelas duas projeções $\operatorname{pr}_i$), perdem-se
irremediavelmente as características particulares do par original
$(p_1,p_2)$. É, portanto, importante encontrar uma generalização da noção de
relação de equivalência que se aplique diretamente ao par definido por um
grupo de operadores de $\mathcal C$.

\subsection*{4. Pré-relações de equivalência}
\addcontentsline{toc}{subsection}{4. Pré-relações de equivalência}

Recordemos que uma categoria na qual todo morfismo é um isomorfismo se chama
\emph{grupoide}. Por outro lado, uma categoria deve ser definida como formada
por dois conjuntos subjacentes $(X,R)$, o conjunto de \emph{objetos} e o
conjunto de \emph{setas}, munidos das seguintes estruturas:

\src{106}
\begin{enumerate}
\item[(i)] Um par de aplicações
\[
  p_1,p_2:R\rightrightarrows X
\]
chamadas \emph{aplicação fonte} e \emph{aplicação alvo},

\item[(ii)] Uma aplicação
\[
  \pi:(R,p_2)\times_X(R,p_1)\longrightarrow R
\]
chamada \emph{aplicação de composição}.
\end{enumerate}
Esses dados devem satisfazer axiomas bem conhecidos, que não repetiremos aqui
e que poderiam ser explicitados pela comutatividade de certos diagramas e pela
existência de uma aplicação $D:X\to R$ (necessariamente única) que torna
comutativos outros dois diagramas. $D$ corresponde à passagem de um objeto à
seta identidade correspondente e satisfaz
\[
  p_1D=p_2D=\operatorname{id}_X.
\]
Dizer que a categoria é um grupoide significa então que existe uma aplicação
\[
  s:R\longrightarrow R
\]
(necessariamente única), chamada a \emph{inversão} de $R$, que transforma cada
seta em uma seta inversa, o que poderia ser explicitado pela comutatividade de
outros quatro diagramas, formados com $s$, $\Delta$ e os dados precedentes, e
dos quais os dois primeiros se escrevem:
\[
  p_1s=p_2,\qquad p_2s=p_1.
\]

Recordados esses pontos, as definições gerais em ([2], A, \S~1) mostram, em
particular, o que se deve entender por uma estrutura de
$\mathcal C$-\emph{categoria}, respectivamente de
$\mathcal C$-\emph{grupoide}, sobre um par de objetos $X,R$ de uma categoria
qualquer $\mathcal C$: trata-se, por definição, do dado, para todo objeto $T$
de $\mathcal C$, de uma estrutura de categoria, respectivamente de grupoide,
no sentido conjuntista, cujo conjunto de objetos é $X(T)$ e o conjunto de
setas é $R(T)$, essas estruturas
\guillemotleft variando funtorialmente\guillemotright{} com $T$ variável.
Isso implica, portanto, a definição de dois morfismos
\[
  p_1,p_2:R\rightrightarrows X
\]
chamados \emph{morfismo fonte} e \emph{morfismo alvo} e, quando existe o
produto fibrado desejado, de um morfismo

\src{107}
\[
  \pi:(R,p_2)\times_X(R,p_1)\longrightarrow R
\]
chamado \emph{morfismo de composição}. Esses três dados bastam então para
determinar a estrutura de categoria sobre $X,R$, sendo o axioma a impor a
esses dados o seguinte: para todo $T$, os três dados correspondentes para
$X(T),R(T)$ definem sobre esse par de conjuntos uma estrutura de categoria
(respectivamente, de grupoide). Quando for o caso, isso pode ser expresso pela
comutatividade de certos diagramas que envolvem um morfismo univocamente
determinado
\[
  D:X\longrightarrow R
\]
e, no caso dos grupoides, um morfismo univocamente determinado
\[
  s:R\longrightarrow R;
\]
diagramas que se explicitam como no caso «conjuntista». Essa interpretação
trabalhosa dos axiomas é, felizmente, inútil na prática. O único interesse
teórico da possibilidade de expressar os dados e os axiomas a partir de
morfismos e igualdades de morfismos entre certos produtos fibrados é o
seguinte: se temos um funtor $F:\mathcal C\to\mathcal C'$ exato à esquerda
(isto é, que comuta com produtos finitos e produtos fibrados), ele transforma
uma $\mathcal C$-categoria (respectivamente, um $\mathcal C$-grupoide) em uma
$\mathcal C'$-categoria (respectivamente, um $\mathcal C'$-grupoide), desde
que existam produtos finitos e produtos fibrados em $\mathcal C$.

Na prática, é importante saber interpretar os morfismos
$p_1,p_2,\pi,D,s$ como \emph{operações simpliciais} em um objeto
semissimplicial ou simplicial adequado de $\mathcal C$, ao menos quando
existem produtos fibrados em $\mathcal C$. Para fixar a terminologia,
introduzamos a \emph{categoria $\mathcal S$ dos simplexos padrão} como a
categoria cujos objetos são os conjuntos finitos da forma
\[
  \Delta_n=[0,n],\qquad\text{onde }n\geq 0
\]
(intervalo dos inteiros de $0$ a $n$), e cujos morfismos ou setas são as
\emph{aplicações arbitrárias} entre esses conjuntos finitos. Observe-se que a
categoria $\mathcal S$ é equivalente à categoria dos conjuntos finitos
\emph{não vazios}, tomando-se como morfismos as aplicações entre conjuntos
finitos. Em $\mathcal S$, existe evidentemente o coproduto de uma família
finita \emph{não vazia} de objetos, bem como o coproduto fibrado de dois
objetos sob um terceiro (operação dual ao produto fibrado). Denotemos por
$\mathcal S'$ a subcategoria de $\mathcal S$ que tem os mesmos objetos, mas
cujos morfismos são as \emph{aplicações crescentes} entre os $\Delta_n$. Essa
categoria é equivalente à categoria dos conjuntos finitos não vazios
totalmente ordenados. Nesta categoria, o

\src{108}
coproduto de dois objetos nunca existe, e o coproduto fibrado de dois objetos
$A,B$ sob um terceiro $C$ não existe em geral (tome-se, por exemplo,
$C=\Delta_0$, $A=B=\Delta_1$, com as duas aplicações estruturais
$u:C\to A$ e $v:C\to B$ iguais). Contudo, ele existe em certos casos; por
exemplo,
\[
  C=\Delta_0,\qquad A=\Delta_m,\qquad B=\Delta_n,\qquad
  u(0)=m,\qquad v(0)=0,
\]
e, nesse caso,
\[
  \Delta_m\amalg_{\Delta_0}\Delta_n=\Delta_{m+n}.
\]

Um \emph{objeto simplicial} (respectivamente, um \emph{objeto
semissimplicial}) em uma categoria $\mathcal C$ é, por definição, um funtor
contravariante $K$ de $\mathcal S$ (respectivamente, de $\mathcal S'$) em
$\mathcal C$. Portanto, um objeto simplicial define por restrição um objeto
semissimplicial, mas a primeira noção difere essencialmente da segunda pela
presença de \emph{operações de simetria} nos $K_n=K(\Delta_n)$,
correspondentes às imagens, pelo funtor $K$, dos elementos do grupo simétrico
de $n+1$ letras (considerado como o grupo de automorfismos de $\Delta_n$ em
$\mathcal S$).

Para todo $n$, seja $\Delta'_n$ (respectivamente, $\Delta''_n$) a categoria
finita cujo conjunto de objetos é $\Delta_n$, definida pela relação de ordem
caótica (respectivamente, pela ordem total natural) em $\Delta_n$ (sendo o
conjunto de setas o grafo dessa relação de ordem). É evidente que $\Delta'_n$
(respectivamente, $\Delta''_n$) depende funtorialmente do objeto $\Delta_n$ de
$\mathcal S$ (respectivamente, de $\mathcal S'$). Se $Z$ for uma categoria,
então $\operatorname{Hom}(\Delta'_n,Z)$ (respectivamente,
$\operatorname{Hom}(\Delta''_n,Z)$), quando $\Delta_n$ varia, é um funtor de
$\mathcal S$ (respectivamente, de $\mathcal S'$) para a categoria dos
conjuntos, isto é, um conjunto simplicial (respectivamente, um conjunto
semissimplicial), dito \emph{associado à categoria $Z$}; denotemo-los por $Z'$
e $Z''$. Há, além disso, um homomorfismo natural evidente do conjunto
semissimplicial associado a $Z'$ para $Z''$, e este último é um isomorfismo se,
e somente se, $Z$ for um grupoide. Posto isso:

\result{PROPOSIÇÃO}{4.1} O funtor $Z\longmapsto Z''$ da categoria das
categorias para a categoria dos conjuntos semissimpliciais é plenamente fiel
e define uma equivalência entre a categoria das \emph{categorias} e a
categoria dos conjuntos semissimpliciais, isto é, dos funtores contravariantes
$K$ de $\mathcal S'$ para $(\mathrm{Set})$ que transformam os coprodutos
fibrados $A\amalg_C B$ do tipo especificado acima em produtos fibrados de
conjuntos. Do mesmo modo, o funtor $Z\longmapsto Z'$ da categoria dos
grupoides para a categoria dos conjuntos simpliciais é plenamente fiel e
define uma equivalência entre a categoria dos \emph{grupoides} e a categoria
dos

\src{109}
conjuntos simpliciais, isto é, funtores contravariantes $K$ de $\mathcal S$
em $(\mathrm{Ens})$, que transformam somas amalgamadas em produtos fibrados.

Pode-se, portanto, considerar as categorias como conjuntos semissimpliciais
particulares, e os grupoides como conjuntos simpliciais particulares, com a
condição, naturalmente, de raciocinar «a menos de isomorfismo», como é de rigor
quando se interpretam certas estruturas em termos de outras. O procedimento
usual de redução ao caso conjuntista implica então:

\result{COROLÁRIO}{4.2} O enunciado precedente continua válido quando se
substituem as categorias, os grupoides e os conjuntos simpliciais por
$\mathcal C$-categorias, $\mathcal C$-grupoides e $\mathcal C$-objetos
simpliciais, contanto que existam produtos fibrados em $\mathcal C$.

O objeto semissimplicial $K$ em $\mathcal C$ associado a uma categoria
$(X,R,\ldots)$ em $\mathcal C$ pode ser explicitado interpretando-se a
componente $K_n=K(\Delta_n)$ de $K$ como a $(n+1)$-ésima potência fibrada de
$(R,p_1)$ sobre $X$, ou melhor, pela fórmula de recorrência
\[
  K_0=R,\qquad
  K_n=(K_{n-1},p_n^{(n-1)})\times_X(R,p_1),
\]
onde os $p_i^{(n-1)}$ ($0\leq i\leq n-1$) são as projeções naturais de
$K_{n-1}$ em $X$ (que também se definem por recorrência). Dessa forma,
$p_1,p_2,\pi,D,s$ são interpretados como as operações simpliciais que
correspondem respectivamente aos morfismos em $\mathcal S$: a face $0$ de
$\Delta_1$, a face $1$ de $\Delta_1$, a face $(0,2)$ de $\Delta_2$, a
degenerescência $\Delta_1\to\Delta_0$ e a simetria de $\Delta_1$. Todas as
demais operações semissimpliciais (respectivamente simpliciais) deduzem-se
formalmente das quatro (respectivamente cinco) precedentes por composição e
produtos fibrados.

Chama-se agora \emph{pré-relação de equivalência} em um objeto $X$ de uma
categoria o dado de um grupoide cujo objeto dos objetos (se assim se pode
dizer) é $X$. Tal dado compreende, portanto, entre outras coisas, um objeto $R$
e dois morfismos
\[
  p_1,p_2:R\rightrightarrows X.
\]
Observe-se, porém, que esses dados, por si sós, não determinam a estrutura
considerada, ao contrário do que ocorre com os pares de equivalência. Nesta
exposição, interessamo-nos por essa noção com o objetivo de obter critérios de
possibilidade de passagem ao quociente, isto é, de formação do conúcleo do par
$(p_1,p_2)$. O enunciado do problema, portanto, não faz uso dos dados
adicionais inerentes a um grupoide. Na demonstração dos resultados que se
seguem,

\src{110}
utilizaremos, contudo, esses dados adicionais e, em particular, as operações
simpliciais (incluindo as operações de simetria) até a dimensão $3$ (nas quais
intervém a quarta potência fibrada de $R$ sobre $X$).

Uma relação de equivalência em um objeto $X$ de $\mathcal C$ define uma
pré-relação de equivalência: basta verificá-lo no caso conjuntista, no qual se
associa a uma relação de equivalência em um conjunto $X$ o grupoide cujo
conjunto de objetos é $X$ e cujo conjunto de setas é o grafo da relação de
equivalência.

Um $\mathcal C$-monoide $G$ que opera sobre um objeto $X$ de $\mathcal C$
define uma $\mathcal C$-categoria cujos objetos subjacentes são $R=G\times X$
e $X$ (desde que $G\times X$ exista), e que é um $\mathcal C$-grupoide se, e
somente se, $G$ for um grupo. Basta, ainda, fazer a verificação no caso
conjuntista. Define-se então a composição das setas $(g,a)$ e
$(g',g\cdot a)$ por
\[
  (g',g\cdot a)(g,a)=(g'g,a),
\]
isto é, se $a,b\in X$, então $\operatorname{Hom}(a,b)$ é, por definição, o
transportador de $a$ em $b$, e os morfismos se compõem mediante a composição
dos elementos de $G$.

\medskip
\noindent\textbf{OBSERVAÇÃO. ---} Podem-se evitar as dificuldades lógicas
levantadas por um enunciado como (4.1), subentendendo-se que todos os objetos
considerados se encontram em um «Universo» fixado (que é ele próprio um
conjunto).

\fsection{5}{Quociente por uma relação de equivalência finita e plana}

\result{TEOREMA}{5.1} Sejam $X=\operatorname{Spec}(B)$ um esquema afim e $R$
uma pré-relação de equivalência em $X$, cuja componente $R_1$ é afim, digamos
$R_1=\operatorname{Spec}(C)$. Suponha-se que a primeira projeção
$p_1:R_1\to X$ seja um morfismo finito e localmente livre; em outras palavras,
que o homomorfismo de anéis correspondente $p_1^*:B\to C$ faça de $C$ um
$B$-módulo projetivo de tipo finito. Seja $A$ o subanel de $B$ que é o núcleo
do par de homomorfismos
\[
  p_1^*,p_2^*:B\rightrightarrows C
\]
(isto é, o conjunto dos elementos $b$ tais que
$p_1^*(b)=p_2^*(b)$). Seja $Y=\operatorname{Spec}(A)$, e seja $f:X\to Y$ o
morfismo definido pela inclusão de $A$ em $B$. Nessas condições:
\begin{enumerate}
\item[(i)] $B$ é integral sobre $A$, isto é, $f$ é um morfismo integral.

\item[(ii)] O morfismo $f$ é sobrejetivo, e suas fibras são as classes de
equivalência

\src{111}
conjuntistas $p_2\allowbreak(p_1^{-1}(x))$ em $X$ módulo $R$, e a topologia
de $Y$ é a topologia quociente da de $X$.

\item[(iii)] $Y$ é o quociente de $X$ por $R$ na categoria dos pré-esquemas.

\item[(iv)] Se $R$ provém de uma \emph{relação de equivalência}, então o
morfismo $f:X\to Y$ é finito e localmente livre (isto é, $B$ é um $A$-módulo
projetivo finitamente gerado), e a relação de equivalência é efetiva, isto é,
$R_1\to X\times_Y X$ é um isomorfismo.
\end{enumerate}

Este teorema generaliza o conhecido teorema relativo ao caso de um grupo
finito $G$ que atua por automorfismos no anel $B$ e ao anel de invariantes
$A$, e sua demonstração é análoga à demonstração conhecida. Pode-se precisar
(iii) da seguinte maneira:

\result{COROLÁRIO}{5.2} O morfismo canônico
$R_1\to X\times_Y X$ é \emph{sobrejetivo}.

Seja ainda $R$ uma pré-relação de equivalência «finita e localmente livre» em
$X$, sendo agora $X$ um pré-esquema arbitrário. Suponhamos que se possa
encontrar um pré-esquema $Y$ e um morfismo $f:X\to Y$ tais que $fp_1=fp_2$ e,
além disso, que seja exata a sequência de homomorfismos de feixes de anéis
sobre $Y$:
\[
  \mathcal O_Y\longrightarrow f_*(\mathcal O_X)
  \rightrightarrows g_*(\mathcal O_R)
\]
(onde $g=fp_i$). Resulta então do teorema que se têm conclusões (i) a (iv)
análogas às do teorema; em particular, por (iii), $Y$ é o quociente de $X$ por
$R$ e, portanto, fica determinado a menos de isomorfismo único. Diz-se, nessas
condições, que a pré-relação de equivalência $R$ em $X$ é \emph{admissível}.
Com essa definição:

\result{TEOREMA}{5.3} Sejam $X$ um pré-esquema e $R$ uma pré-relação de
equivalência em $X$ tal que $p_1:R_1\to X$ seja um morfismo finito e localmente
livre. Para que $R$ seja admissível, é necessário e suficiente que toda classe
de equivalência conjuntista $p_2(p_1^{-1}(x))$ em $X$ módulo $R$ esteja
contida em um aberto afim (condição sempre verificada se todo subconjunto
finito de $X$ estiver contido em um aberto afim, por exemplo, se $X$ for
quase-projetivo sobre um esquema afim).

Mostra-se, com efeito, sem dificuldade que toda classe de equivalência módulo
$R$ em $X$ está então contida em um aberto afim \emph{estável por $R$}, e
constrói-se o quociente $Y$ pela colagem das partes obtidas com o auxílio do
Teorema 5.1.

\src{112}

\result{COROLÁRIO}{5.4} Suponhamos satisfeita essa condição e que, além disso,
$R$ provenha de uma relação de equivalência. Então esta última é efetiva, isto
é, $R_1\to X\times_Y X$ é um isomorfismo, e $f:X\to Y$ é um morfismo finito
localmente livre.

Conclui-se imediatamente, por «descida»:

\result{COROLÁRIO}{5.5} Sob as condições de (5.4), $X$ tem posto $n$ em todo
ponto sobre $Y$ se e somente se $(R_1,p_1)$ tem posto $n$ em todo ponto sobre
$X$. Se $X,R_1$ forem $Z$-pré-esquemas e $p_1,p_2$ forem $Z$-morfismos, sendo
portanto $Y$ um $Z$-pré-esquema, então $X$ é plano sobre $Z$ se e somente se
$Y$ o for.

Em resumo:

\result{ESCÓLIO}{} O dado de um morfismo finito, localmente livre e
sobrejetivo $f:X\to Y$ de pré-esquemas equivale ao dado de um pré-esquema $X$
munido de uma relação de equivalência $R$ tal que $p_1:R\to X$ seja finito e
localmente livre, e que toda classe $p_2(p_1^{-1}(x))$ esteja contida em um
aberto afim.

\medskip
\noindent\textbf{OBSERVAÇÕES 5.6.}

\smallskip
\noindent $1^\circ$ Não foi necessário fazer nenhuma hipótese Noetheriana.

\smallskip
\noindent $2^\circ$ Esta noção de passagem ao quociente inclui, como caso
particular, a «descida inseparável» de CARTIER, que corresponde à determinação
dos morfismos finitos localmente livres $f:X\to Y$ tais que
$f_*(\mathcal O_X)$ admita uma $p$-base em relação a $\mathcal O_Y$ (sendo $X$
um pré-esquema dado cujo feixe $\mathcal O_X$ é anulado pelo número primo
$p>0$). Observe-se que o resultado se exprime facilmente sem hipótese de
regularidade sobre os anéis locais e sem supor que $X$ seja um esquema
algébrico sobre um corpo. A teoria de JACOBSON--BOURBAKI é obtida tomando-se
como $X$ o espectro de um corpo de característica $p$.

\smallskip
\noindent $3^\circ$ GABRIEL havia obtido anteriormente um caso particular do
teorema (5.3) na teoria da passagem ao quociente para grupos comutativos
finitos sobre um corpo $k$. [Compare-se (7.4).]

\fsection{6}{Quociente por uma relação de equivalência própria e plana}

\result{TEOREMA}{6.1} Sejam $S$ um pré-esquema localmente Noetheriano, $X$
um $S$-esquema quasi-projetivo, $R$ uma pré-relação de equivalência em $X$,
tal que:

\src{113}
\begin{enumerate}
\item[(a)] $p_1:R_1\to X$ é próprio e plano;
\item[(b)] $R_1\to X\times_S X$ é um morfismo finito ou, o que equivale em
virtude de (a), tem fibras finitas (condição automaticamente satisfeita se
$R$ provém de uma relação de equivalência).
\end{enumerate}
Nessas condições:
\begin{enumerate}
\item[(i)] $X/R=Y$ existe e (se $S$ é Noetheriano e se $R$ provém de uma
relação de equivalência) é quase-projetivo sobre $S$.\footnote{O texto impresso
não contém a segunda hipótese. A errata vinculada à exposição 212 esclarece
que a quase-projetividade parece então demonstrada apenas quando $R$ provém
de uma relação de equivalência.}

\item[(ii)] O morfismo canônico $f:X\to Y$ é sobrejetivo, próprio e aberto;
suas fibras são as classes de equivalência $p_2(p_1^{-1}(x))$ em $X$ módulo
$R$; portanto, $Y$ identifica-se com o espaço topológico quociente de $X$ pela
relação de equivalência conjuntista definida por $R$. Por fim,
$R_1\to X\times_Y X$ é sobrejetivo.

\item[(iii)] Se $R$ provém de uma relação de equivalência, esta última é
efetiva, isto é, $R_1\to X\times_Y X$ é um isomorfismo e, além disso,
$f:X\to Y$ é plano (portanto, fielmente plano).
\end{enumerate}

Para a demonstração, reduz-se a (5.1) por meio de quase-seções adequadas de $X$
para $R$, sendo a demonstração análoga à construção dos grupos algébricos
quocientes no Seminário Chevalley.

Em resumo:

\result{ESCÓLIO}{} Seja $X$ quase-projetivo sobre $S$ localmente Noetheriano. O
dado de um morfismo próprio, fielmente plano e sobrejetivo $f:X\to Y$ de $X$
para um $S$-pré-esquema $Y$ equivale ao dado de uma relação de equivalência $R$
em $X$, tal que $p_1:R\to X$ seja próprio e plano.

O mesmo método fornece o resultado seguinte:

\result{TEOREMA}{6.2} Sejam $S$ um pré-esquema Noetheriano, $X$ um
pré-esquema de tipo finito sobre $S$ e $R$ uma pré-relação de equivalência no
$S$-pré-esquema $X$; supõe-se que:
\begin{enumerate}
\item[(a)] $p_1:R_1\to X$ é plano e de tipo finito;
\item[(b)] o morfismo $R_1\to X\times_S X$ é quase-finito (isto é, tem fibras
finitas).
\end{enumerate}
Então existe um aberto $U$ denso em $X$, \emph{saturado} para $R$, tal que:
\begin{enumerate}
\item[(i)] Se $R_U$ é a pré-relação de equivalência induzida por $R$ em $U$,
então $U/R_U$ existe e é de tipo finito sobre $S$.

\item[(ii)] O morfismo canônico $U\to U/R_U$ é sobrejetivo e aberto; suas
fibras são as classes de equivalência conjuntistas para $R_U$ (portanto,
$U/R_U$ é o espaço topológico quociente de $U$ pela relação de equivalência
conjuntista definida por $R_U$); por fim, o morfismo
\[
  (R_U)_1\longrightarrow U\times_{U/R_U}U
\]
é sobrejetivo.

\src{114}

\item[(iii)] Se $R$ provém de uma relação de equivalência, pode-se supor que
$U\to U/R_U$ seja fielmente plano e que $R_U$ seja efetiva.
\end{enumerate}

Este é um resultado de natureza essencialmente «birracional».

\medskip
\noindent\textbf{OBSERVAÇÕES 6.3.}

\smallskip
\noindent $1^\circ$ Ignoro se, em (6.1) e (6.2), a hipótese (b) é inútil. Na
passagem ao quociente por grupos, ela nos obriga praticamente a nos
restringirmos ao caso em que os estabilizadores são todos grupos finitos.

\smallskip
\noindent $2^\circ$ Pode-se perguntar se não há resultados análogos a (6.1) e
(6.2) sem hipótese de planicidade. Não conheço nenhum contraexemplo nessa
direção. Em contrapartida, mesmo conservando a hipótese de planicidade e
restringindo-nos a relações de equivalência tais que $p_1:R\to X$ seja plano e
quase-finito (mas não finito), sendo $X$ afim, pode ocorrer que $R$ não seja
efetiva: tomem-se as relações de equivalência induzidas nos abertos afins que
recobrem a variedade de Nagata (na qual um grupo de ordem dois opera de modo
«não admissível»).

\fsection{7}{Aplicações}

Como dissemos na introdução, a aplicação mais importante de (6.1) é a
construção de esquemas de Picard, bem como de soluções para vários outros
problemas de moduli, aos quais voltaremos posteriormente.

Obtém-se uma demonstração simples do seguinte resultado de SHIMURA:

\result{PROPOSIÇÃO}{7.1} Seja $A$ um esquema abeliano definido sobre um anel de
valorização discreta $V$, com corpo de frações $K$. Então todo esquema
abeliano $B'$ sobre $K$, isógeno a um quociente de $A\otimes_V K$, «tem boa
redução sobre $V$», isto é, é isomorfo a algum $B\otimes_V A$,
\footnote{O texto impresso traz $B\otimes_V A$. Como $B$ é um $V$-esquema e a
afirmação diz respeito à sua fibra genérica, o fator de base deve ser $K$.} onde $B$ é um esquema
abeliano sobre $V$ (essencialmente único, recordemo-lo).

Pode-se supor que $B'$ seja o quociente de $A_K$ por um subesquema em grupos
$C'$. (N.B.: $C'$ não será, em geral, «reduzido», isto é, seus anéis locais
terão elementos nilpotentes.) Considere-se o subesquema fechado $C$ de $A$ que
é o «fecho» de $C'$, isto é, o menor subesquema fechado de $A$ tal que $C_K$
contenha $C'$. Teremos $C_K=C'$, e do fato de que $V$ é um anel de valorização
discreta deduz-se facilmente que $C$ é um subesquema em grupos de $A$ sobre
$V$. Como $A$ é próprio sobre $\operatorname{Spec}(V)=S$, o mesmo vale para
$C$; por outro lado, $A$ é até projetivo sobre $S$. Pode-se então aplicar
(6.1) para construir $A/C=B$,

\src{115}
que é o $B$ procurado.

Por fim, argumentos essencialmente conhecidos permitem deduzir de (6.2) o
seguinte resultado:

\result{TEOREMA}{7.2} Sejam $S$ o espectro de um anel Artiniano, $F$ e $G$
esquemas em grupos de tipo finito sobre $S$, e $u:F\to G$ um homomorfismo de
esquemas em grupos sobre $S$. Suponhamos:
\begin{enumerate}
\item[(i)] $F$ é plano sobre $S$;
\item[(ii)] o núcleo de $u$ é finito.
\end{enumerate}
Nessas condições, existe o esquema quociente $G/F$; o morfismo canônico
$G\to G/F$ é sobrejetivo e aberto, e suas fibras são as classes de equivalência
conjuntistas definidas pelas ações à direita de $F$ sobre $G$. Por fim, se $u$
é um monomorfismo, então $G\to G/F$ é plano e o morfismo
\[
  G\times F\longrightarrow G\times_{G/F}G
\]
é um isomorfismo; em outras palavras, $G$ é um espaço homogêneo principal
sobre $G/F$, com grupo estrutural (atuando à direita) $F$, ou melhor,
$F\times_S(G/F)$ considerado como esquema em grupos sobre $G/F$ (cf. [1], B,
\S~6).

\result{COROLÁRIO}{7.3} Nessas condições, para que $G$ seja plano sobre $S$,
é necessário e suficiente que $G/F$ o seja. Se essa condição for satisfeita,
a passagem ao quociente por $F$ comuta com toda extensão da base $S$ e, se
$F$ for um subgrupo normal de $G$, $G/F$ poderá ser munido de uma estrutura
de \emph{grupo quociente} de $G$ por $F$.

A situação é particularmente simples quando $S$ é o espectro de um corpo,
pois então todo $S$-pré-esquema é automaticamente plano sobre $S$. Obtém-se:

\result{COROLÁRIO}{7.4} Sejam $F,G$ esquemas em grupos de tipo finito sobre um
corpo $k$, e seja $u:F\to G$ um homomorfismo de $k$-grupos. Então $u$ se
fatora como $F\to F'\to G$, onde $F\to F'$ é o homomorfismo de passagem ao
quociente pelo subgrupo fechado $\operatorname{Ker}u$ de $F$, e $F'\to G$ é
um homomorfismo de grupos que é uma imersão fechada. Existe o quociente
$G/F=G/F'$. O formalismo usual (do tipo dos teoremas de Noether) é válido
entre os grupos algébricos sobre $k$.

Este resultado permite tratar de maneira uniforme a passagem ao quociente nos
grupos algébricos no sentido clássico (isto é, os que são irredutíveis sobre
$k$ e suaves sobre $k$) e a passagem ao quociente por subgrupos
«infinitesimais»

\src{116}
considerados por CARTIER. É vantajoso considerar as «hiperálgebras»
introduzidas por CARTIER, na esteira dos trabalhos de DIEUDONNÉ sobre os grupos
formais, como grupos na categoria dos esquemas formais sobre $k$ e, quando for
o caso (se corresponderem a hiperálgebras de posto finito sobre $k$), como
grupos algébricos finitos sobre $k$.

\fsection{8}{Uma conjectura}

Ela responde à necessidade de saber passar ao quociente pelo grupo projetivo
que atua em certos subesquemas dos «esquemas de Hilbert» (estes últimos
substituindo, na teoria dos esquemas, as variedades de Chow).

Sejam $S$ um pré-esquema e $n$ um inteiro. A todo pré-esquema $S'$ sobre $S$,
façamos corresponder o grupo
$\operatorname{Gl}(n,\Gamma(S',\mathcal O_{S'}))$ das matrizes invertíveis de
$n$ linhas e $n$ colunas com entradas no anel das seções de
$\mathcal O_{S'}$. Obtém-se assim um funtor contravariante em $S'$, que se
mostra facilmente representável; ele corresponde, portanto, a um esquema em
grupos sobre $S$, de fato afim sobre $S$, denotado por $\operatorname{Gl}(n)_S$.
Sua formação é compatível com a mudança de base, de modo que, na realidade,
tudo provém de um esquema em grupos sobre $\mathbf Z$, denotado por
$\operatorname{Gl}(n)$. O grupo $\operatorname{Gl}(1)$, chamado grupo
multiplicativo e frequentemente denotado por $\mathbf G_m$, corresponde ao
funtor $S\longmapsto\Gamma(S,\mathcal O_S)^*$, grupo das «unidades» sobre $S$.
Há um homomorfismo evidente
$\operatorname{Gl}(1)\to\operatorname{Gl}(n)$, e constrói-se facilmente o
grupo quociente, denotado por $\operatorname{GP}(n-1)$ e chamado
\emph{grupo projetivo de grau $n-1$ sobre $\mathbf Z$}. Ele representa o
funtor que associa a $S$ o grupo das seções do feixe
$\widetilde{\operatorname{Gl}(n)}_S/
\widetilde{\operatorname{Gl}(1)}_S$, onde
$\widetilde{\operatorname{Gl}(n)}_S$ denota o feixe dos germes de seções de
$\operatorname{Gl}(n)_S$ sobre $S$. (É preciso observar que as seções de
$\operatorname{GP}(n-1)_S$ sobre $S$ não provêm, em geral, de seções de
$\operatorname{Gl}(n)_S$ sobre $S$!) Note-se que se pode provar que
$\operatorname{GP}(n-1)$ representa também o funtor
$S\longmapsto\operatorname{Aut}_S(\mathbf P_S^{n-1})$, onde
$\mathbf P_S^{n-1}$ é o espaço projetivo padrão de dimensão relativa $n-1$
sobre $S$, ao menos quando $S$ é Noetheriano. É dessa maneira que ele intervém
na teoria de módulos.

Seja $S$ um esquema Noetheriano, que se pode, se desejado, supor afim, e seja
$X$ um $S$-pré-esquema quase-projetivo, munido de um feixe invertível
$\mathcal L$ muito amplo relativamente a $S$. Suponhamos que o grupo
$G=\operatorname{GP}(n)_S$ atue em $X$ e, ao mesmo tempo, em $\mathcal L$ (de
modo compatível com sua ação em $X$), e que atue \emph{livremente} em
$S$.\footnote{O texto impresso diz «em $S$». A ação considerada e a condição
de quociente dizem respeito a $X$.}

\src{117}

\result{CONJECTURA}{8.1}\footnote{A errata vinculada à exposição 212 assinala
que esta conjectura é «decididamente falsa». Indica também que o fato positivo
da observação 8.1 foi desde então demonstrado e remete aos trabalhos de Mumford
e a seus refinamentos.} Nas condições precedentes:

\smallskip
\noindent $1^\circ$ \emph{A relação de equivalência definida por $G$ é
efetiva, o quociente $Y=X/G$ é de tipo finito sobre $S$ e o morfismo canônico
$f:X\to Y$ é plano e sobrejetivo} (portanto, $X$ torna-se um fibrado principal
homogêneo sobre $Y$, de grupo
$G\times_S Y=\operatorname{GP}(n)_Y$).

\smallskip
\noindent $2^\circ$ Seja $\mathcal L'$ o feixe invertível sobre $Y$ deduzido
de $\mathcal L$ por «descida fielmente plana» ao longo de $f$ (cf. [1], B,
\S~th. 1). Então $\mathcal L'$ é \emph{pré-amplo sobre $Y$ relativamente a
$S$}, isto é, existem um inteiro $m$ e um morfismo quase-finito de $Y$ para um
esquema projetivo padrão conveniente $\mathbf P_S^N$, tais que
$(\mathcal L')^{\otimes m}$ seja isomorfo à imagem inversa de
$\mathcal O_{\mathbf P_S^N}(1)$.

Observe-se que, mesmo se $X$ é separado sobre $S$, pode ocorrer que $Y$ não
seja separado sobre $S$ (situação que ocorre em problemas de módulos nada
patológicos). Se $1^\circ$ se verifica, então $Y$ é separado se e somente se a
relação de equivalência definida por $G$ tem gráfico fechado, isto é, se
$G\times X\to X\times X$ tem imagem fechada (trata-se então de uma imersão
fechada). Se $Y$ é separado, então $\mathcal L'$ é pré-amplo sobre $Y$
relativamente a $S$ se e somente se é amplo, isto é, se uma potência tensorial
adequada define uma imersão projetiva. Nos problemas de módulos mencionados na
introdução, pode-se mostrar que a relação de equivalência obtida tem de fato um
gráfico fechado.

\medskip
\noindent\textbf{OBSERVAÇÕES 8.1. ---} Supusemos
$G=\operatorname{GP}(n)_S$ para fixar as ideias e porque este é atualmente o
caso mais importante na prática. A hipótese razoável a fazer sobre $G$
pareceria antes que $G$ fosse uma das «formas» sobre $S$ de um dos grupos de
Tohokû (cuja construção sobre os inteiros também foi realizada por CHEVALLEY).
O único fato positivo que conheço na direção da conjectura precedente é o
seguinte: seja $X$ \emph{um esquema afim sobre um corpo $k$ de característica
$0$, no qual o grupo $\operatorname{Gl}(n)_k$ ou
$\operatorname{GP}(n-1)_k$ atua livremente. Então a relação de equivalência
definida por $G$ é efetiva, o quociente $X/G$ também é afim e o morfismo
$X\to X/G$ é plano e sobrejetivo}. A demonstração utiliza o fato seguinte (que
por enquanto só está demonstrado em característica $0$): quando se faz $G$
atuar sobre o anel de coordenadas afim $A$ de $G$, considerado como espaço
vetorial sobre $k$, a representação trivial de $G$ ocorre apenas uma vez (em
uma série de composição de um subespaço vetorial de dimensão finita sobre $k$,
estável sob $G$). Parece possível que um uso sistemático da

\src{118}
a teoria das representações lineares de $G$ acabará por fornecer uma
demonstração da conjectura, ao menos quando se está sobre um corpo de base.
Quando já não se está sobre um corpo de base, o conferencista nada sabe.

\section*{BIBLIOGRAFIA}
\addcontentsline{toc}{section}{Bibliografia}

\begin{description}
\item[{[1]}] GROTHENDIECK (Alexander). --- Techniques de descente et théorèmes
d'existence en géométrie algébrique, I: Généralités, Séminaire Bourbaki,
1959/60, n\textsuperscript{o} 190, 29 p.

\item[{[2]}] GROTHENDIECK (Alexander). --- Techniques de descente et théorèmes
d'existence en géométrie algébrique, II: Le théorème d'existence en théorie
formelle des modules, Séminaire Bourbaki, 1959/60,
n\textsuperscript{o} 195, 22 p.

\item[{[3]}] GROTHENDIECK (A.) e DIEUDONNÉ (J.). --- Éléments de géométrie
algébrique, I: Le langage des schémas. --- Paris, Presses universitaires de
France, 1960 (Institut des hautes Études scientifiques, Publications
mathématiques, 4).
\end{description}

